\chapter{路网层次结构特性与双曲空间建模理论分析}

本章对路网层次结构的多尺度特性进行系统分析，论证双曲空间与路网层次结构的适配性，并提出隐式捕获与显式约束两种层次建模范式。这些分析为第四章和第五章的方法设计提供统一的理论根基和动机论证。

\section{路网层次结构特性分析}

真实路网中存在从个体路段到宏观功能区的多尺度层次组织。最底层是路段（Segment），即独立的道路单元，承载具体的通行功能，数量通常从数千到数万不等。中间层是局部区域（Locality），由地理上相邻的路段自然聚合而成的结构性子区域——例如一个立交桥群、一个小区内部路网单元，或一组围绕同一交叉口的路段集合。最顶层是功能区（Region），覆盖大片地理空间的功能性分区，如商业区、居民区或工业区，通常与城市规划中的土地利用类型对应。

这三个层次之间存在明确的包含关系：每个功能区包含若干局部区域，每个局部区域包含若干路段。这种"Region $\supseteq$ Locality $\supseteq$ Segment"的包含结构构成了路网的天然层次体系。在语义上，这一包含关系对应着"通用-具体"的概念梯度——功能区代表最通用的地理概念（数量少，通常 $|\mathcal{R}| \approx 30$），局部区域代表中等粒度的结构概念（$|\mathcal{C}| \approx 300$），路段代表最具体的空间单元（$|\mathcal{S}|$ 从数千到数万）。

路网的层次组织呈现出树状生长特性。从少数动脉级主干道（高速公路、快速路）出发，道路系统逐级分叉：主干道连接次级干道，次级干道连接集散路，集散路连接居民区支路。在这种分支过程中，每一级的道路数量大致按指数规律增长。以一个典型的中等规模城市为例，高速公路和快速路通常仅有数条，主干道有数十条，而居民区支路可达数千条。

Barth\'{e}lemy\cite{Barthelemy2011} 的研究从空间网络理论角度证实了这一特性。城市道路系统的生长过程本质上是一种层次化的空间填充过程：主干道首先划定城市骨架，随后在骨架之间逐级填充更密集的次级道路。这一过程产生的拓扑结构，其分支模式与生物系统中的树状脉管系统具有深刻的类比关系。

路网的层次结构体现在两个相对独立的维度上。结构层次（Structural Hierarchy）由路网的拓扑连接决定，反映的是路段之间的空间邻接和连通关系，可以通过谱聚类等图划分方法得到。功能层次（Functional Hierarchy）由交通功能决定，反映的是道路在交通系统中承担的角色（如主干道、支路），通常需要轨迹数据或道路类型标注来确定。

结构层次具有较高的稳定性，因为路网的物理拓扑一旦建成便不易改变。功能层次则具有一定的动态性——同一路段在早高峰和深夜可能承担不同的交通功能。然而，在日均或周均的时间尺度上，功能层次也呈现出相当的稳定性。本文的两项工作均假设路网为静态图，在这一假设下，两种层次结构均可被视为稳定的先验知识。

\section{双曲几何理论基础}

双曲空间是具有恒定负曲率的黎曼流形。与欧氏空间的平坦性（曲率为零）和球面空间的正曲率不同，双曲空间的负曲率赋予其独特的几何性质。最显著的性质是体积的指数增长：在 $n$ 维双曲空间中，以某点为中心、半径为 $r$ 的测地球的体积满足
\begin{equation}
    V_{\mathbb{H}}(r) \propto e^{(n-1)r\sqrt{|\kappa|}}
\end{equation}
其中 $\kappa < 0$ 为曲率。这与欧氏空间中 $V_{\mathbb{E}}(r) \propto r^n$ 的多项式增长形成鲜明对比。

双曲空间的另一个重要性质是三角形的角度和小于 $\pi$。这意味着双曲空间中的"距离感"与欧氏空间不同：远离中心的区域具有更大的"有效面积"，这为在低维空间中嵌入大量叶节点提供了几何基础。

庞加莱球模型定义 $d$ 维双曲空间为开球 $\mathbb{B}_c^d := \{\mathbf{x} \in \mathbb{R}^d : c\|\mathbf{x}\|^2 < 1\}$，其中 $c \geq 0$ 为曲率参数。庞加莱球满足共形性，即球内曲线的角度与欧氏测量一致，但距离向边界急剧拉伸。

庞加莱球中两点 $\mathbf{x}, \mathbf{y} \in \mathbb{B}_c^d$ 之间的测地距离为：
\begin{equation}
    d_c(\mathbf{x}, \mathbf{y}) = \frac{2}{\sqrt{c}} \tanh^{-1}\left(\sqrt{c}\|\mathbf{-x} \oplus_c \mathbf{y}\|\right)
\end{equation}
其中 $\oplus_c$ 为莫比乌斯加法（M\"{o}bius Addition），定义为：
\begin{equation}
    \mathbf{x} \oplus_c \mathbf{y} := \frac{(1 + 2c\langle \mathbf{x},\mathbf{y}\rangle + c\|\mathbf{y}\|^2)\mathbf{x} + (1 - c\|\mathbf{x}\|^2)\mathbf{y}}{1 + 2c\langle \mathbf{x},\mathbf{y}\rangle + c^2\|\mathbf{x}\|^2\|\mathbf{y}\|^2}
\end{equation}

莫比乌斯标量乘法定义为：
\begin{equation}
    r \otimes_c \mathbf{x} := \frac{1}{\sqrt{c}}\tanh\left(r\tanh^{-1}(\sqrt{c}\|\mathbf{x}\|)\right)\frac{\mathbf{x}}{\|\mathbf{x}\|}
\end{equation}

指数映射（从原点处切空间到庞加莱球）和对数映射（从庞加莱球到原点处切空间）分别为：
\begin{equation}
    \exp_\mathbf{0}^c(\mathbf{v}) = \tanh\left(\sqrt{c}\|\mathbf{v}\|\right)\frac{\mathbf{v}}{\sqrt{c}\|\mathbf{v}\|}, \quad
    \log_\mathbf{0}^c(\mathbf{x}) = \frac{1}{\sqrt{c}}\tanh^{-1}(\sqrt{c}\|\mathbf{x}\|)\frac{\mathbf{x}}{\|\mathbf{x}\|}
\end{equation}

Lorentz 模型 $\mathbb{L}^n$ 定义为 $(n+1)$ 维 Minkowski 时空中的上双曲面：
\begin{equation}
    \mathbb{L}^n = \left\{\mathbf{p} \in \mathbb{R}^{n+1} : \langle \mathbf{p}, \mathbf{p} \rangle_{\mathbb{L}} = -\frac{1}{\kappa},\ p_0 > 0\right\}
\end{equation}
其中 Lorentz 内积定义为 $\langle \mathbf{p}, \mathbf{q} \rangle_{\mathbb{L}} = -p_0q_0 + \sum_{i=1}^n p_iq_i$，$\kappa > 0$ 为曲率。两点之间的 Lorentz 距离为：
\begin{equation}
    d_{\mathbb{L}}(\mathbf{p}, \mathbf{q}) = \sqrt{\frac{1}{\kappa}} \cdot \cosh^{-1}(-\kappa\langle \mathbf{p}, \mathbf{q} \rangle_{\mathbb{L}})
\end{equation}

Lorentz 模型中的指数映射为：
\begin{equation}
    \exp_{\mathbf{p}}^\kappa(\mathbf{v}) = \cosh(\sqrt{\kappa}\|\mathbf{v}\|_{\mathbb{L}})\mathbf{p} + \frac{\sinh(\sqrt{\kappa}\|\mathbf{v}\|_{\mathbb{L}})}{\sqrt{\kappa}\|\mathbf{v}\|_{\mathbb{L}}}\mathbf{v}
\end{equation}

庞加莱球模型和 Lorentz 模型在数学上等价，但在数值计算中表现出不同的稳定性特征。庞加莱球模型的计算在球内进行，当嵌入点接近单位球边界（$\|\mathbf{x}\| \to 1/\sqrt{c}$）时，莫比乌斯运算涉及的分母趋近于零，容易导致数值溢出。这一问题在需要精确表示深层次叶节点（即靠近边界的点）时尤为突出。

Lorentz 模型的指数映射和对数映射具有封闭形式，数值稳定性更优。Peng 等\cite{peng2022hyperbolic} 的系统比较表明，在深层次结构嵌入任务中，Lorentz 模型的训练稳定性显著优于庞加莱球模型。

基于这一对比，本文在两项工作中做出了不同的模型选择。第四章的 HRGNN 采用庞加莱球模型，因为其任务（道路类型分类）更关注判别性特征学习而非深层次组织，且庞加莱球的共形性有助于注意力权重的计算。第五章的 HHRoad 采用 Lorentz 模型，因为其需要在三层层次结构中精确组织嵌入的几何关系，对数值稳定性要求更高。

\section{路网层次结构与双曲空间的适配性论证}

路网层次结构的核心特征是节点数量随层次深度的指数增长。在路网的三层体系中，功能区数量 $|\mathcal{R}|$ 最少（$\sim$30），局部区域数量 $|\mathcal{C}|$ 居中（$\sim$300），路段数量 $|\mathcal{S}|$ 最多（数千至数万）。从功能区到路段，数量增长大致遵循指数规律。

双曲空间的指数体积增长恰好与这一特性匹配。在双曲空间中，距原点半径 $r$ 处的"可用空间"随 $r$ 指数增长，这意味着大量路段嵌入可以在双曲空间的外围区域以较低的嵌入失真分布，同时与位于靠近原点处的少数功能区嵌入保持合理的层次距离关系。相比之下，欧氏空间中半径 $r$ 处的可用空间仅随 $r^d$ 多项式增长，要容纳同等数量的嵌入点需要更高的维度或接受更大的失真。

Sala 等\cite{Sala2018} 的理论分析为这一适配性提供了严格的量化依据：双曲嵌入能以 $O(d)$ 维度以任意精度近似嵌入 $n$ 节点的树，而欧氏嵌入需要 $\Omega(\log n / d)$ 的失真下界。

$\delta$-双曲性（$\delta$-Hyperbolicity）\cite{Borassi2015} 是衡量图拓扑结构与树的相似程度的数学工具。其定义基于四点条件：对任意四点 $x, y, z, w$，图中最大的 $\delta$ 值满足
\begin{equation}
    (x|z)_w \geq \min\{(x|y)_w, (y|z)_w\} - \delta
\end{equation}
其中 $(x|y)_w = \frac{1}{2}(d(x,w) + d(y,w) - d(x,y))$ 为 Gromov 积。$\delta$ 值越小，图越接近树结构（$\delta=0$ 即为完美的树），越适合双曲嵌入。

\begin{figure}[t]
\centering
\begin{tabular}{@{}c@{\hspace{0.5em}}c@{}}
\includegraphics[width=0.40\textwidth]{figures/figure2a} &
\includegraphics[width=0.40\textwidth]{figures/figure2b} \\
(a) & (b)
\end{tabular}
\caption{全球 $\delta$-双曲性统计分析。(a) 大洲分布分析及整体分布叠加（灰色箱体），揭示了显著的形态差异。(b) 散点图考察网络规模（节点数量，单位千）与 $\delta$-双曲性之间的关系，颜色代表不同大洲。}
\label{fig:hyperbolicity}
\end{figure}


为验证路网的树状特性，本文对 Urbanity 全球城市网络数据集\cite{Yap2023} 中 29 个国家 50 个城市进行了大规模 $\delta$-双曲性分析。图3.1分析结果揭示了三个重要发现。

第一，路网的 $\delta$ 值普遍较低。威尼斯的 $\delta=3.0$，慕尼黑的 $\delta=17.5$，雅典的 $\delta=32.0$，涵盖了从高度双曲到低双曲的广泛范围，但即使是 $\delta$ 较大的现代规划城市，其路网仍然具有显著的层次特征。

第二，$\delta$-双曲性与城市形态密切相关，但与网络规模无强相关性。中世纪有机生长的城市（如威尼斯）具有更低的 $\delta$ 值和更强的树状特性，而现代规划的网格状城市（如雅典）$\delta$ 值较高。城市规模（路段数量）并不决定 $\delta$ 值，这表明路网的几何特性独立于网络规模，由城市形态和规划模式决定。

第三，即使是 $\delta$ 值较大的路网，其层次结构仍然存在。$\delta$ 值较大意味着路网偏离完美树的程度较大，但不意味着层次结构消失——网格状路网仍然具有功能分层（高速公路、干道、支路）和结构聚类特性。

在双曲空间中，到原点的距离天然编码层次深度。以 Lorentz 模型为例，原点（North Pole）$\mathbf{0} = (\sqrt{1/\kappa}, 0, \ldots, 0)^\top$ 位于双曲面的"顶部"，距原点越远的点位于层次结构的越深层。具体地，如果功能区嵌入于距原点较近处（$d_{\mathbb{L}}(\mathbf{h}_r, \mathbf{0})$ 较小），局部区域嵌入于中间距离，路段嵌入于远离原点处，那么层次深度就通过嵌入点到原点的距离得到了自然编码。

这种映射关系不需要额外的约束即可在优化过程中自然出现，前提是优化目标中包含对层次信息的监督信号。庞加莱球模型中也存在类似的映射：靠近球心的点对应层次结构中的一般性概念（如功能区），靠近球边界的点对应特定性概念（如路段）。

将层次结构嵌入的根本问题可以形式化为以下问题：给定一棵 $n$ 节点的树 $T$，将其节点映射为 $d$ 维嵌入空间中的点，使得嵌入后的距离尽可能保持树中的最短路径距离。嵌入的失真（Distortion）定义为所有节点对的距离比的最大值。

在欧氏空间 $\mathbb{R}^d$ 中，Bourgain 定理保证存在 $O(\log^2 n)$ 失真的嵌入，但 Linial 等证明了失真下界为 $\Omega(\log n / d)$——即使维度 $d$ 增大，失真仍至少以 $\log n$ 的速度增长。这意味着对于大规模路网（$n$ 可达数万），欧氏嵌入的失真不可避免地较大。

在双曲空间 $\mathbb{H}^d$ 中，Sarkar\cite{sarkar2011low} 证明了仅用 2 维双曲空间即可以 $(1+\epsilon)$ 的失真嵌入任意树，代价是坐标精度为 $O(n \log n / \epsilon)$ 位。这一理论结果表明，双曲空间在表达层次结构方面具有根本性的几何优势——它能够在低维空间中以近乎零失真嵌入树状结构，这是欧氏空间无法做到的。

对于路网这类不完全是树但具有强树状特性的图，双曲嵌入虽然无法达到零失真，但其失真仍然显著低于欧氏嵌入，且这一优势随路网规模增大而更加明显。

\section{路网层次关系的双曲建模范式探讨}

基于上述理论分析，本文提出路网层次结构的双曲建模可以沿两种范式展开。

第一种是\textbf{隐式捕获范式}。这一范式利用双曲空间的度量性质——双曲距离天然反映层次深度差异——来隐式编码层次信息。具体地，在双曲空间中进行消息传递和注意力计算时，双曲距离作为节点间相似性的度量，自然地使得层次位置相近的节点获得更大的注意力权重。层次结构的编码并非通过专门的约束实现，而是作为双曲空间几何特性的"副产品"。

第二种是\textbf{显式约束范式}。这一范式在隐式利用双曲空间度量性质的基础上，进一步通过蕴含锥等几何工具对层次包含关系施加明确的数学约束。"功能区包含局部区域"的语义关系被转化为严格的几何不等式——子概念嵌入必须位于父概念的蕴含锥内。层次结构的编码不仅依赖空间的背景几何，还通过优化目标中的硬约束直接实现。

隐式范式的优势在于简洁和通用。由于不需要预定义层次划分或设计复杂的层次约束，隐式范式适合于层次信息作为辅助特征而非核心建模目标的场景。例如，在道路类型分类任务中，模型的主要目标是学习路段的判别性特征，层次结构信息起辅助作用——隐式范式通过双曲空间的几何特性自然融入层次先验，无需额外的层次划分预处理和蕴含损失设计。

隐式范式的局限在于层次关系的表达力度有限。双曲距离可以反映层次深度差异，但无法表达"包含"这一偏序关系的方向性——两个双曲距离相近的点可能是同层次的兄弟节点，也可能是不同层次的父子节点。这种歧义在需要精确层次组织的任务（如需要从功能区语义向路段传递层次信息的通用表示学习任务）中会成为瓶颈。

显式范式的优势在于层次关系的表达精确且可解释。蕴含锥机制将"父-子"包含关系转化为可验证的几何条件——子概念嵌入是否位于父概念蕴含锥内可以通过角度比较直接判定。这种精确性使得模型的层次组织可以通过几何度量直接解读，为后续分析和决策提供了可解释的依据。

显式范式的代价在于需要预定义层次划分（哪些路段属于哪个局部区域、哪个功能区）以及设计相应的蕴含损失函数。这意味着模型需要更多的先验信息输入和更复杂的训练目标。显式范式适合于层次结构建模是核心目标的场景，如需要产出可迁移至多种下游任务的通用路段嵌入。

两种范式之间存在自然的递进关系。隐式范式验证了双曲空间对路网层次建模的基本有效性——如果将路段嵌入从欧氏空间移至双曲空间就能带来性能提升，那么双曲空间的几何特性确实与路网层次结构匹配。在这一验证的基础上，显式范式进一步追问：能否通过设计专门的几何约束，将双曲空间的层次建模潜力从"被动利用"提升到"主动约束"？

这一递进逻辑构成了本文两项工作的组织脉络：第四章的 HRGNN 采用隐式范式，验证双曲空间的有效性并揭示其局限；第五章的 HHRoad 采用显式范式，通过蕴含锥机制克服隐式建模的不足，实现对路网层次关系的精确几何约束。

\section{本章小结}

本章从理论角度系统分析了路网层次结构特性与双曲空间的适配性。路网呈现出从路段到局部区域到功能区的多尺度层次组织，节点数量随层级深度呈指数增长，拓扑结构具有可量化的树状特性。双曲空间的指数体积增长、层次深度的距离编码和蕴含锥的偏序建模能力，使其成为路网层次建模的天然几何框架。基于这一分析，本文提出了隐式捕获与显式约束两种建模范式，分别对应第四章和第五章的方法设计，两者形成连贯的研究脉络。